
This chapter reviews the theoretical and empirical literature underpinning Uzbekistan's integration into global value chains (GVCs), with a specific focus on Northeast Asian foreign direct investment (FDI) and the structural constraints of landlocked industrialisation. The review synthesises existing scholarship, identifies contradictions in the application of developmental state theory to Central Asia, and highlights the methodological limitations of current empirical approaches to trade decomposition.

\subsection{Global Value Chain Theory}
Global Value Chain (GVC) theory provides the dominant framework for understanding modern industrialisation. Early formulations by Gereffi (1994) distinguished between producer-driven and buyer-driven commodity chains, which later evolved into a sophisticated typology of governance structures: market, modular, relational, captive, and hierarchy. For emerging economies, the transition from captive to relational governance is the critical threshold for technological upgrading. However, recent scholarship critiques the assumption that insertion into GVCs automatically yields upgrading, noting that lead firms often actively block functional upgrading by developing country suppliers.

\subsection{Developmental State Theory and Industrial Policy}
The East Asian 'miracle' was heavily predicated on the developmental state model, characterized by bureaucratic autonomy, selective industrial policy, and export discipline. While scholars have attempted to apply this framework to Central Asia, the comparison often fails due to the absence of embedded autonomy and the prevalence of rent-seeking over export-oriented industrialisation. The resurgence of industrial policy in the post-Washington Consensus era necessitates a re-evaluation of state intervention in Uzbekistan's post-2017 reform context.

\subsection{FDI and Technology Spillovers}
The endogenous growth literature posits that FDI is a primary conduit for technology transfer and human capital accumulation. However, empirical evidence on horizontal and vertical spillovers is highly mixed. Studies suggest that positive spillovers only materialize when the host economy possesses sufficient 'absorptive capacity'—threshold levels of human capital and institutional readiness. In the context of transition economies, enclave FDI—where foreign firms operate in isolated special economic zones without domestic backward linkages—remains a persistent risk.

\subsection{Export-Led vs Import-Led Industrialisation}
Traditional development economics emphasizes export-led industrialisation as the only viable path to sustained growth. Conversely, import-substituting industrialisation (ISI) is generally viewed as an exhausted paradigm. Yet, contemporary literature highlights a hybrid model: 'import-led exports', where access to high-quality intermediate goods is a prerequisite for export competitiveness. For landlocked countries, the tariff penalty on imported intermediates disproportionately depresses export sophistication.

\subsection{Vietnam and Thailand GVC Integration}
Vietnam and Thailand represent the standard benchmarks for successful GVC embedding in developing Asia. Thailand's Eastern Seaboard development utilized targeted Japanese FDI to build an automotive cluster. Vietnam's post-Doi Moi trajectory leveraged Samsung's relocation to dominate electronics assembly. The literature emphasizes that both countries paired aggressive trade liberalisation with massive investments in maritime logistics and SEZ infrastructure, a structural advantage unavailable to Central Asia.

\subsection{Northeast Asian Production Networks}
Northeast Asia (China, Japan, and South Korea) functions as a highly integrated trilateral production network, often conceptualized through the 'Flying Geese' model. As wages rise in China, production is relocating to peripheral nodes (the 'China+1' strategy). However, empirical studies indicate that Chinese outward FDI differs fundamentally from Japanese and Korean FDI: the latter is heavily manufacturing-focused, while Chinese capital often flows into natural resources and infrastructure (the Belt and Road Initiative). \subsection{Central Asian Structural Constraints}
The literature on landlocked developing countries (LLDCs) consistently identifies the 'distance penalty' as the primary barrier to industrialisation. Transit costs function as an exogenous tariff, isolating LLDCs from time-sensitive, fragmented global production networks. For Uzbekistan, the double-landlocked geographic condition means that traditional maritime-based export models are structurally invalid, necessitating alternative corridors such as the Middle Corridor (TITR) or air-freight industrialisation.

\subsection{SEZ and Industrial Cluster Literature}
Special Economic Zones (SEZs) are a ubiquitous policy tool, yet the empirical consensus on their effectiveness is deeply polarized. Successful SEZs function as institutional test-beds and overcome infrastructure bottlenecks, whereas failed SEZs operate as mere tax havens with minimal linkage to the domestic economy. In the post-Soviet context, SEZs have historically underperformed due to poor site selection, regulatory unpredictability, and failure to attract anchor tenants.

\subsection{Research Gap}
The existing literature exhibits three critical gaps. First, while extensive scholarship analyzes Southeast Asian GVC integration, the mechanisms of LLDC integration into Northeast Asian networks remain under-theorized. Second, empirical studies on Uzbekistan overwhelmingly focus on macroeconomic liberalisation rather than disaggregated, firm-level GVC governance structures. Third, comparative benchmarking between Central Asian and ASEAN industrial models is largely absent from the literature. This dissertation addresses these gaps by quantifying the structural asymmetries in Uzbekistan's trilateral engagement.
